\documentclass[10pt]{article}
\usepackage[margin=20mm]{geometry}
\usepackage{amsmath,amssymb,fontspec,unicode-math,xurl,hyperref}
\IfFontExistsTF{DejaVu Serif}{\setmainfont{DejaVu Serif}}{\setmainfont{Latin Modern Roman}}
\setmathfont{latinmodern-math.otf}
\setlength{\parindent}{0pt}\setlength{\parskip}{6pt}\sloppy
\title{Evaluation and representation checks for the few-shot image detector}
\author{AI-assisted technical review}
\date{}
\begin{document}
\maketitle
Abstract: The prototype computation, squared-distance classifier and episode indexing agree with the stated method in independent finite checks. The headline gain is 11.6 percentage points under the paper’s evaluation protocol. Two release-level inconsistencies remain: testing omits the described resize, and visualization removes the trained projection used for classification. Neither establishes failure of historical accuracy results. The episodic support budget and unequal baseline training information require careful interpretation. Scientific adjudication remains pending.\par
\section*{Identity and scope}
This report reviews Shiyu Wu, Jing Liu, Jing Li and Yequan Wang, Few-Shot Learner Generalizes Across AI-Generated Image Detection, ICML 2025, PMLR 267. The technical body on pages 1–8 of the 12-page proceedings PDF was read, together with the available public source at https://github.com/teheperinko541/Few-Shot-AIGI-Detector, commit `b545c05f3c927ef67c1b00f9a8badf3b68c5f4b3`. The remaining PDF pages contain acknowledgments, impact and references; there is no technical appendix in this version. Static inspection covered the prototype utility, train/test drivers, dataset loaders, argument parser, model/save utilities, train/evaluation/rebuild shell scripts, requirements, README and visualization driver. Logger internals were not reviewed. No repository program, model, checkpoint, image dataset, GPU training or t-SNE computation was run. Original Figures 2–6 and Algorithm 1 on source pages 3, 4, 7 and 8 were actually viewed.\par
\section*{Checks supporting the local method}
For feature vector \(z\) and real/fake prototype means \(c_R,c_F\), the two-class squared-distance softmax has log-odds \[\log\frac{p(F\mid z)}{p(R\mid z)}
 =\|z-c_R\|^2-\|z-c_F\|^2
 =2z^\top(c_F-c_R)+\|c_R\|^2-\|c_F\|^2.\] Thus recomputing prototypes adapts an affine decision boundary in the learned feature space without a gradient update. It does not imply a universal separation guarantee for arbitrary new generators. The sample mean minimizes the sum of squared distances, consistent with equation (1) and the selected distance in Section 4.2.\par
An independently written NumPy check uses real supports \((-2,-1)\), fake supports \((1,2)\) and queries \((-0.4,0.4)\). The centers are \(-1.5,1.5\); squared distances are \((1.21,3.61)\) and \((3.61,1.21)\). The correct-class probabilities are both 0.916827, and the affine logit differences are \(-2.4,2.4\), agreeing with direct distances. An independent identity-index example also checks the train driver’s class-major flattening followed by task/sample/class reshaping. The resulting query labels alternate in the same order as its repeated class labels, and the first-\(K\) support IDs do not overlap the remaining query IDs. The release’s evaluation mode disables training-mode batch-statistic updates. No within-episode support/query self-inclusion is established. Actual file duplicates and dataset splits were not audited.\par
The six Table 1 entries average to 84.0667\% for 10-shot FSD, 77.0500\% for zero-shot FSD, and 72.4833\% for LARE2. They agree with the rounded 84.1, 77.1 and 72.5. The reported 11.6 improvement is in percentage points, corresponding to a relative increase of 16.0\% over 72.5; the additional few-shot gain over the reported zero-shot FSD value is 7.0 points. These calculations validate table arithmetic, not independent model accuracy.\par
\section*{Problem 1: the released test transform omits the stated resize}
Severity: Minor; a supported reproduction discrepancy with unmeasured historical effect. Section 4.2, page 5, says that test images are resized and then center-cropped to \(224\times224\). In `datasets/\_\_init\_\_.py`, the validation/test loader instead applies only `CenterCrop(224)` and `ToTensor()`; `test.py` directly uses this loader on each class’s validation directory. The training loader has a resize, so this is not merely a shared omitted description of preprocessing.\par
For a \(1024\times1024\) input, the release’s native \(224\times224\) crop retains 4.7852\% of the original area. Resizing to \(256\times256\) before the same crop covers approximately \(896\times896\) original coordinates, or 76.5625\%, a sixteenfold geometric area difference. This is a coordinate-footprint calculation, not a claim that all benchmark images have that resolution or that either transformed image was run through the detector. Spatial scale and retained image content are different inputs to the detector, so these pipelines are not interchangeable in general.\par
If every supplied file was already resized to the required size by an unrecorded preprocessing step, the discrepancy could disappear. The included rebuild script creates symlinks and does not perform that resize, and the README does not specify such a step. The historical dataset/preprocessing chain is therefore needed. A reproduction should explicitly choose the documented resize-and-crop path or revise the method description to match the intended native crop, then rerun the affected comparison. No performance repair has been tested.\par
\section*{Problem 2: visualization uses a different representation}
Severity: Minor; release-level mismatch in the evidence for the learned metric space. Figure 2 specifies 1024-dimensional features, and the train/test drivers instantiate a ResNet-50 with a 1024-output head. Figure 4 and Section 4.5 describe visualizing the final-layer features of the model. However, `visualization/tsne\_script.py` loads that model and then replaces `model.fc` with `torch.nn.Identity()` before feature extraction. This removes the trained projection used by the actual prototype classifier, producing backbone features instead.\par
A learned projection need not preserve neighbors: take backbone points \(a=(0,100)\), \(b=(1,0)\), \(c=(2,100)\) and projection \(W=(1,0)\). Then \(\|a-c\|=2<\|a-b\|\simeq100.005\), but \(|Wa-Wb|=1<|Wa-Wc|=2\). The independent calculation verifies this reversal. Thus a visualization of pre-head features cannot directly certify the geometry of the post-head metric representation. This does not prove that the displayed historical figures used the current script, or that their qualitative clusters would disappear. The repair is to retain the classifier’s projection when studying its metric space, or accurately label the plotted features as backbone features and provide a separate post-head analysis. Neither actual visualization was regenerated here. Even correctly sourced t-SNE plots are qualitative two-dimensional evidence, not quantitative guarantees of generalization.\par
\section*{Protocol clarifications and unresolved comparisons}
The source and Algorithm 1 use \(K\) support images per class. A binary 10-shot episode therefore uses ten fake and ten real supports; only the ten fake images may be new acquisitions if the real pool is already available. In the released evaluator, each new shuffled batch receives a fresh support subset, and all query predictions are then aggregated. Across \(m\) such episodes, up to \(mK\) distinct labeled fake images are used, not one fixed global set of \(K\). An independent two-episode ID example with \(K=2\) consumes four distinct fake supports without any query overlap. This is valid episodic few-shot evaluation, but does not by itself measure deployment with ten total fake labels reused for an entire generator’s test stream. The fixed 1:3 support/query ratio in Section 4.6 is compatible with this episodic interpretation. A fixed-support evaluation and uncertainty over repeated support draws would establish the stronger acquisition-budget claim.\par
Other concrete release settings also need an explicit run mapping. The supplied evaluation shell script does not override parser defaults of five supports and fifteen queries, hence it requests 5-shot evaluation. It passes a seed, but `test.py` never calls the seed-setting helper; a static syntax-tree inspection confirms that absence. The test loader shuffles. The training example requests 50,000 steps instead of the paper’s 200,000, though it can be configured to that value. The README/train comment describes real images from SD v1.4/v1.5, whereas the paper says the real subsets are aggregated. None of these defaults proves which configuration produced a published checkpoint. The zero-shot training-prototype procedure is described in the paper but not implemented by the inspected test driver.\par
Section 4.3 reports FSD trained jointly on the non-test generator classes, while its baseline comparison averages separate models trained on individual non-test subsets. FSD also receives labeled target-generator supports in the few-shot condition; the paper explicitly acknowledges that difference. Consequently the 11.6-point comparison combines a different training information regime with target support adaptation. It does not isolate the causal effect of the prototype method. A joint-data baseline, equally budgeted target adaptation and quantitative multiclass versus binary ablation would resolve that attribution. No matched rerun is available, so this is Needs Evidence, not a claim that FSD would lose its advantage.\par
The release requests a multiclass average-precision metric with ten thresholds; the paper also states a coarse threshold step. The exact library version, averaging convention and output scores were not available for reproducing AP. No exact rank-based AP equivalence is assumed, and no incorrect historical AP value is established.\par
\section*{Disposition}
Assessment | Repair verification | Author adoption | Remaining work\par
P1 Minor | Geometry difference checked | Unknown | Historical transforms; rerun\par
P2 Minor | Projection counterexample checked | Unknown | Exact feature provenance\par
Protocol: Needs Evidence | No empirical repair | Unknown | Fixed-support, matched baselines\par
The local prototype algebra and disclosed accuracy arithmetic survive the checks. No Major failure of the detection method has been demonstrated. The two concrete release discrepancies and remaining protocol questions limit reproduction and interpretation. Scientific adjudication remains pending, and no human review is recorded.\par
\section*{Evidence and review limits}
The source is the frozen PMLR proceedings PDF, identified by SHA-256\par
> `a81a2651ec22c3a4de05dcf3c78db67b4269957efd7c16d1ea1990f3d785e889`\par
Physical PDF page numbers are one based and coincide with the paper’s printed page numbers in this version. The check program was written for this review; no author program was executed. Published training runs, data collection and reported performance were not reproduced. It is not independent acceptance. No human review or author adoption is recorded. Scientific adjudication remains pending.\par
9 Shiyu Wu; Jing Liu; Jing Li; Yequan Wang. Few-Shot Learner Generalizes Across AI-Generated Image Detection. Proceedings of the 42nd International Conference on Machine Learning, PMLR 267 (2025). Published proceedings PDF, frozen for this review on 25 October 2026. https://proceedings.mlr.press/v267/wu25r.html.\par
\end{document}
